\section{Introduction}
\label{sec:introduction}

Variable binding is the recurring technical obstacle in any
formalization touching the lambda calculus: essentially every
representation choice (named terms modulo alpha equivalence,
de Bruijn indices, locally nameless terms, higher-order abstract
syntax, nominal sets, \ldots) trades one set of proof obligations for
another, and results proved for one representation rarely transfer
automatically to the others. This has led to a large but fragmented
body of formalization work, often redone from scratch for each new
proof assistant or each new project, with little of it packaged as a
reusable library.

CSLib is an effort to build such a reusable, foundational library for
computer science results in the Lean~4 theorem prover, in the spirit
of Mathlib for mathematics~\cite{demoura2026a}. This paper reports on
two contributions to CSLib's treatment of the untyped lambda calculus:

\begin{enumerate}
  \item a machine-checked formalization of Crole's result that five
    common literature definitions of alpha equivalence on named lambda terms
    are equivalent~\cite{crole2012a}, which to our knowledge
    had not previously been formalized in a proof assistant; and
  \item an isomorphism, following Norrish's proof
    pearl~\cite{norrish2007a}, between the quotient of named
    terms by a chosen alpha equivalence and a de Bruijn
    representation, connecting the named and nameless ends of
    CSLib's binding infrastructure.
\end{enumerate}

We present both results as part of a longer-term plan: to grow CSLib
into an idiomatic, well-reviewed, and maintained foundation that
other researchers can build on directly, rather than reformalizing
lambda calculus basics for each new project. Section~\ref{sec:future-work}
sketches this plan in more detail, including further equivalences in
the style of Crole's for SK combinator calculus and other binding
representations, and a longer-term goal of formalizing a substantial
part of Barendregt's reference treatment of the untyped lambda
calculus~\cite{barendregt1985a}.

\paragraph{Outline.}
Section~\ref{sec:cslib} introduces CSLib and its current named
representation of the untyped lambda calculus.
Section~\ref{sec:alpha-equivalences} presents the formalization of
Crole's alpha-equivalence equalities.
Section~\ref{sec:de-bruijn} presents the isomorphism with the de
Bruijn representation. Section~\ref{sec:related-work} discusses
related formalization efforts, Section~\ref{sec:future-work} lays
out planned extensions, and Section~\ref{sec:conclusion} concludes.
